% --- CLIP background: needed before Stable Diffusion's text conditioning ---
\begin{frame}{Background: CLIP (Contrastive Language--Image Pre-training)}
    \textbf{Why we need it:} Stable Diffusion does not learn ``text'' from scratch --- it
    relies on a pre-trained \textbf{CLIP text encoder} to turn a prompt into a vector the
    UNet can condition on. So before the pipeline, we need to know what CLIP is.

    \vspace{0.6em}
    \begin{columns}[T,onlytextwidth]
        \column{0.56\textwidth}
            \textbf{What CLIP does}
            \begin{itemize}\itemsep0.25em
                \item Trains an \textbf{image encoder} and a \textbf{text encoder} jointly on
                      $\sim$400M (image, caption) pairs scraped from the web.
                \item \textbf{Contrastive objective:} pull the embeddings of a matching
                      image--caption pair together, push mismatched pairs apart.
                \item Result: a \textbf{shared embedding space} where a picture of a cat and the
                      text ``a photo of a cat'' land close together.
            \end{itemize}
        \column{0.44\textwidth}
            \textbf{Why it matters for diffusion}
            \begin{itemize}\itemsep0.25em
                \item Gives a \textbf{semantic, language-aligned} representation of the prompt.
                \item That text embedding is injected into the denoiser via
                      \textbf{cross-attention}, steering generation toward the prompt.
                \item Enables \textbf{zero-shot}, open-vocabulary text-to-image control.
            \end{itemize}
            \vspace{0.4em}
            {\footnotesize Radford et al., ``Learning Transferable Visual Models From Natural
            Language Supervision,'' ICML 2021 (\href{https://arxiv.org/abs/2103.00020}{arXiv:2103.00020}).}
    \end{columns}
\end{frame}
